%% Complete proofs for the migrated propositional syntax/semantics results.

\begin{annexitem}{annex:p2.1.1}{thm:p02-parenthesis-counting}{Parenthesis counting}
\begin{proof}
Structural induction is immediate for variables, \(\bottom\), and negation.
For \((\varphi\circ\psi)\), the outer pair preserves the balance of the two
subformulas.  At the main \(\circ\), the outer left parenthesis precedes the
balanced string \(\varphi\), giving excess one.  At an occurrence within
either subformula, that outer parenthesis adds one to the positive prefix
excess supplied by induction (the preceding \(\varphi\) is balanced in the
right case).  Hence both claims hold.
\end{proof}
\end{annexitem}

\begin{annexitem}{annex:p2.1.2}{thm:p02-unique-readability}{Unique readability}
\begin{proof}
The four clauses have disjoint initial tokens: a variable, \(\bottom\),
\(\neg\), or a left parenthesis.  Removing \(\neg\) fixes the unary
subformula.  In a binary formula, scan with parenthesis depth including the
outer left parenthesis.  The main connective occurs at depth one because its
left subformula is balanced.  Annex~\ref{annex:p2.1.1} puts every connective
internal to either immediate subformula at depth greater than one.  Thus the
unique depth-one connective fixes the split, the connective, and both
subformulas.
\end{proof}
\end{annexitem}

\begin{annexitem}{annex:p2.1.3}{thm:p02-propositional-coincidence}{Coincidence for propositional formulas}
\begin{proof}
Induct structurally: the hypothesis handles variables, \(\bottom\) is fixed,
and connective cases follow from their induction hypotheses and truth clauses.
\end{proof}
\end{annexitem}

\begin{annexitem}{annex:p2.1.4}{thm:p02-clause-satisfaction}{Satisfaction of clauses and clause conjunctions}
\begin{proof}
If \(C=\varnothing\), then
\(\operatorname{form}_{\mathsf{CPL}}(C)=\bottom\), and both sides of the
first equivalence are false.  Otherwise write the increasing-code listing as
\(L_1,\ldots,L_k\).  Downward induction gives
\[
 v\entails D_i\quad\text{iff}\quad
 v\entails L_j\text{ for some }j\in\{i,\ldots,k\}.
\]
The base is \(D_k=L_k\); the step is the \(\lor\)-clause.  Taking \(i=1\)
proves the first claim.  A second downward induction along
\(N_i=\operatorname{form}_{\mathsf{CPL}}(C_i)\land N_{i+1}\), using the
\(\land\)-clause, proves that \(N_1\) is true exactly when every clause
formula is true.
\end{proof}
\end{annexitem}

\begin{annexitem}{annex:p2.1.5}{thm:p02-negation-normal-form}{Negation normal form}
\begin{proof}
Replace each \(\alpha\to\beta\) with \(\neg\alpha\lor\beta\).  Induct
structurally, removing double negations and applying De Morgan to push others
inward.  Truth clauses verify every rewrite; each lowers the number of
connectives below a negation.  Termination leaves only literals and units,
hence an equivalent NNF.
\end{proof}
\end{annexitem}

\begin{annexitem}{annex:p2.1.6}{thm:p02-conjunctive-disjunctive-normal-form}{Conjunctive and disjunctive normal forms}
\begin{proof}
Let the variables occurring in \(\varphi\) be
\(p_{i_1},\ldots,p_{i_n}\), where \(i_1<\cdots<i_n\); the indices need not
form an initial segment.  For an assignment \(s\) on this finite set put
\[
 \ell_j^s=
 \begin{cases}p_{i_j},&s(p_{i_j})=1,\\
 \neg p_{i_j},&s(p_{i_j})=0,
 \end{cases}
 \qquad
 \delta_s=\bigwedge_{j=1}^{n}\ell_j^s.
\]
The finite-connective convention gives \(\delta_s=\topv\) when \(n=0\).
By its truth clauses, a total valuation satisfies \(\delta_s\) exactly when
its restriction to the displayed variables is \(s\).  Order all assignments
lexicographically, and disjoin \(\delta_s\) over precisely those \(s\) that
satisfy \(\varphi\).  The empty disjunction is \(\bottom\).  Coincidence now
shows that this DNF has exactly the valuations of \(\varphi\).

For CNF, for every assignment \(s\) falsifying \(\varphi\), let
\[
 K_s=\{\overline{\ell_1^s},\ldots,
         \overline{\ell_n^s}\},
\]
The repaired ordering detail is essential: because
\(i_j<i_{j+1}\), every code on \(p_{i_j}\) is at most \(2i_j+1\), which
is strictly below every code on \(p_{i_{j+1}}\), at least \(2i_{j+1}\).
Thus the displayed literal order is exactly the increasing-code order used by
\(\operatorname{form}_{\mathsf{CPL}}\).  Annex~\ref{annex:p2.1.4} gives that
this clause formula is false exactly when the valuation restricts to \(s\).
Conjoin these clause formulas over all falsifying assignments, in
lexicographic order.  If there are none, the conjunction is \(\topv\); if
\(n=0\) and \(\varphi\) is false, the sole assignment yields
\(K_s=\varnothing\) and hence the one-clause CNF \(\bottom\).  In every case,
the clause-conjunction lemma and coincidence show that the CNF is satisfied
exactly when \(\varphi\) is.  This preserves R-011: increasing code fixes each
clause, Annex~\ref{annex:p2.1.4} supplies both downward inductions, and the
cases cover validity and both variable-free values.
\end{proof}
\end{annexitem}

\begin{annexitem}{annex:p2.1.7}{thm:p02-functional-basis}{A nullary-inclusive functionally complete basis}
\begin{proof}
For \(f\), take the DNF of its true inputs, replacing disjunction by
\(\neg(\neg\alpha\land\neg\beta)\).  Only \(\neg,\land,\bottom\) remain;
the empty disjunction is \(\bottom\), and \(\neg\bottom\) represents the
nullary true function.
\end{proof}
\end{annexitem}

\begin{annexitem}{annex:p2.1.8}{thm:p02-propositional-duality}{Duality theorem}
\begin{proof}
Induct on the displayed NNF presentation \(\alpha\).  If \(\alpha\) is a
positive literal \(p\), then
\(v\entails p\) iff \(\bar v\not\entails p\).  If it is \(\neg p\), then
\(v\entails\neg p\) iff \(v(p)=0\), iff
\(\bar v(p)=1\), iff \(\bar v\not\entails\neg p\).  The two unit
cases follow from \(\bottom^{\mathrm d}=\topv\) and
\(\topv^{\mathrm d}=\bottom\).  For conjunction,
\[
 v\entails(\beta\land\gamma)^{\mathrm d}
 \quad\text{iff}\quad
 v\entails\beta^{\mathrm d}\text{ or }v\entails\gamma^{\mathrm d}
 \quad\text{iff}\quad
 \bar v\not\entails\beta\land\gamma,
\]
by the induction hypotheses; the disjunction case is dual.  This proves the
displayed semantic identity.

The recursive clauses also give \((\alpha^{\mathrm d})^{\mathrm d}=\alpha\)
by structural induction.  If \(\alpha\Leftrightarrow\beta\), apply the
semantic identity to \(\bar v\) for each \(v\) to obtain
\(\alpha^{\mathrm d}\Leftrightarrow\beta^{\mathrm d}\); the converse follows
from involution.  If \(\alpha\) is valid, the identity makes
\(\alpha^{\mathrm d}\) false under every \(v\).  Conversely, if the dual is
unsatisfiable, apply the identity with \(v=\bar u\) to show that every \(u\)
satisfies \(\alpha\).  Finally, the syntactic clauses swap every outer
conjunction with disjunction, every clause-disjunction with a conjunction of
the same literals, and the two empty units; hence they exchange CNF and DNF.
\end{proof}
\end{annexitem}
